\chapter{Manual de uso}

Este manual describe el despliegue y el funcionamiento del cuadro de mando analítico para la prevención del abandono escolar. Está dirigido a dos perfiles: el administrador del sistema, que debe levantar la infraestructura y poblar los datos, y el técnico de la administración educativa, que utilizará el \textit{dashboard} para la toma de decisiones.

\section{Despliegue del sistema}

\subsection{Requisitos previos}

\begin{itemize}
    \item Tener instalado Docker y Docker Compose.
    \item Clonar el repositorio disponible en \href{https://github.com/andrespazparedes/tfg}{github.com/andrespazparedes/tfg}\footnote{\url{https://github.com/andrespazparedes/tfg}}.
    \item Crear un archivo .env en la raíz del proyecto a partir de la plantilla .env.example. Este archivo contiene las credenciales de las bases de datosy los puertos.
\end{itemize}

\subsection{Levantar los contenedores}

Ejecutar en la terminal, dentro de la carpeta del proyecto:

\begin{verbatim}
docker-compose up -d --build
\end{verbatim}

Este comando levanta los seis contenedores que componen el sistema:

\begin{enumerate}
    \item Base de datos origen (PostgreSQL) con los datos transaccionales simulados.
    \item Data Warehouse (PostgreSQL) con el modelo en estrella.
    \item Apache Hop Web, para ejecutar los procesos ETL.
    \item Backend FastAPI, que expone la API REST.
    \item Frontend React, con la interfaz de usuario.
    \item Proxy Nginx, que unifica el acceso a la aplicación.
\end{enumerate}

Una vez levantados, el \textit{frontend} estará accesible en http://localhost y Apache Hop en http://localhost:8080.

\subsection{Poblar el Datamart}

Los datos no se cargan automáticamente. Es necesario ejecutar los flujos de Apache Hop para extraer, transformar y cargar la información desde la base de datos origen al Data Warehouse.

Para ello, acceder a la interfaz web de Hop (http://localhost:8080), abrir el proyecto y ejecutar el workflow principal, llamado main.hwf. Este workflow orquesta la carga de todas las dimensiones y las tablas de hechos.

En modo incremental, el sistema solo procesa los datos del curso indicado y recupera el histórico acumulado de años anteriores.

\subsection{Personalización}

Todo el código fuente está disponible en el repositorio. Los administradores pueden modificar las reglas de negocio (ponderaciones del riesgo, umbrales de alerta, etc.) editando los \textit{pipeline}s de Hop o el código del \textit{backend}, y reconstruir los contenedores con los cambios.

\section{Uso del cuadro de mando}

\subsection{Acceso y autenticación}

Abrir el navegador en la URL donde se sirve la aplicación (por defecto, http://localhost). La pantalla de inicio solicita correo electrónico y contraseña. Tras autenticarse, la sesión se mantiene mediante un token JWT almacenado en el navegador, por lo que al cerrar y volver a abrir el navegador no es necesario volver a introducir las credenciales si el token sigue siendo válido.

\subsection{Barra lateral y filtros globales}

La barra lateral izquierda permite navegar entre las dos vistas principales (Macro y Micro) y, para los administradores, acceder a la gestión de usuarios. Además, contiene los filtros globales que condicionan todos los datos mostrados:

\begin{itemize}
    \item Curso académico: selecciona el año de análisis.
    \item Tipo de centro: público, concertado o privado.
    \item Ciclo educativo: Primaria, ESO, etc.
    \item Centros concretos: permite elegir uno o varios centros para análisis focalizado.
\end{itemize}

Al modificar cualquier filtro, todas las visualizaciones se actualizan automáticamente.

\subsection{Vista Macro: identificación de centros prioritarios}

Esta vista ofrece una panorámica autonómica y está pensada para responder a la pregunta: ¿qué centros requieren atención prioritaria?

\subsubsection{Indicadores de cabecera}

Tres tarjetas resumen el estado general:
\begin{itemize}
    \item Total de centros que cumplen los filtros.
    \item Número de centros críticos (índice de riesgo igual o superior a 7 sobre 10).
    \item Riesgo medio autonómico.
\end{itemize}

\subsubsection{Panel de alertas institucionales}

Diez tarjetas muestran, para cada métrica de riesgo, cuántos centros superan la mediana autonómica en el número de alumnos afectados. Las métricas son:

\begin{itemize}
    \item Alto riesgo de abandono.
    \item Alumnado con suspensos.
    \item Desfase académico por edad.
    \item Repetidores.
    \item Riesgo socioeconómico alto.
    \item Brecha digital.
    \item Bajo nivel de estudios familiar.
    \item Adaptaciones curriculares.
    \item Repetidores en 1.º y 2.º de primaria.
    \item Suspensos en 1.º y 2.º de primaria.
\end{itemize}

Al pulsar sobre cualquiera de estas tarjetas, el ranking de centros se filtra para mostrar solo aquellos que superan la mediana en esa métrica, facilitando la focalización en problemas concretos.

\subsubsection{Mapa geográfico}

Cada centro aparece en el mapa con un marcador de color (verde, amarillo o rojo) según su nivel de riesgo. Al hacer clic en un marcador se puede acceder a un resumen y navegar directamente a la vista Micro de ese centro.

\subsubsection{Gráficos de análisis macro}

\begin{itemize}
    \item Evolución del índice de riesgo: tendencia del riesgo medio en los últimos cinco años.
    \item Distribución por tipo de centro: barras que comparan centros públicos, concertados y privados.
    \item Matriz de asignación de recursos: gráfico de dispersión que cruza el índice de riesgo (eje X) con el número de alumnos (eje Y). Los centros que combinan alto riesgo y gran volumen de alumnos son los de mayor prioridad.
\end{itemize}

\subsubsection{Ranking de centros}

Lista completa de centros ordenados de mayor a menor índice de riesgo. Cada tarjeta muestra el nombre, tipo, número de alumnos, índice de riesgo y variación respecto al curso anterior. Al pulsar sobre una tarjeta se realiza un \textit{drill-down} a la vista Micro con ese centro preseleccionado.

\subsection{Vista Micro: diagnóstico causal de un centro}

Esta vista responde a la pregunta: ¿por qué este centro tiene el riesgo que tiene? Se accede desde el menú o mediante \textit{drill-down} desde la Macro.

\subsubsection{Indicadores de alumnado}

Diez tarjetas KPI muestran, a escala de alumno, el número de estudiantes afectados por cada factor de riesgo en el curso actual (T0), junto con la variación porcentual respecto al curso anterior (T-1). Esto permite evaluar la evolución temporal de cada problema.

Al pulsar sobre una tarjeta, la tabla de alumnos se filtra para mostrar únicamente los estudiantes que presentan esa alerta, y la vista se desplaza automáticamente hasta la tabla.

\subsubsection{Gráficos de análisis de alumnado}

Ocho gráficos desglosan la composición del riesgo:

\begin{itemize}
    \item Evolución histórica del riesgo medio en los últimos cinco cursos.
    \item Distribución de suspensos: histograma con el número de alumnos que acumulan 0, 1, 2, 3 o más suspensos.
    \item Distribución del nivel de riesgo: gráfico circular con los porcentajes de alumnos en riesgo bajo, medio y alto.
    \item Riesgo por ciclo educativo: comparativa del porcentaje de alumnos en riesgo crítico por ciclo (Primaria, ESO, etc.).
    \item Distribución del nivel de renta: tramos de renta familiar (Muy baja, Baja, Media/Alta).
    \item Correlación renta vs. suspensos: diagrama de dispersión que evidencia la relación entre vulnerabilidad económica y rendimiento.
    \item Brecha digital vs. aprobados: relación entre acceso a dispositivos e internet y tasa de aprobado.
    \item Nivel educativo familiar: distribución del máximo nivel de estudios de los progenitores.
\end{itemize}

Los cuatro últimos gráficos son clave para diferenciar si el problema es principalmente académico (suspensos, repetición) o socioeconómico (renta baja, brecha digital, bajo nivel educativo), lo que orienta la decisión sobre qué tipo de recursos asignar.

\subsubsection{Directorio de alumnos}

Tabla paginada que lista los alumnos que cumplen los filtros activos. Muestra expediente, centro, ciclo, tasa de aprobado, puntuación de riesgo (con código de colores) y badges de alertas activas. Permite ordenar por cualquier columna y filtrar por tipo de alerta mediante un selector rápido. Un contador indica el total de alumnos coincidentes.

\subsection{Gestión de usuarios (solo administradores)}

Accesible únicamente para cuentas con rol de administrador, permite crear, editar y eliminar cuentas de usuario. La cuenta con identificador 1 (superadministrador) está protegida: no puede ser eliminada ni degradada a rol de visualizador, garantizando que siempre haya al menos una cuenta con permisos totales.

\subsection{Flujo de trabajo recomendado}

\begin{enumerate}
    \item El técnico consulta la vista Macro, observa el panel de alertas y el ranking, e identifica los centros que más destacan por riesgo alto o por métricas concretas (por ejemplo, brecha digital o repetidores en primaria).
    \item Mediante \textit{drill-down}, accede a la vista Micro de esos centros.
    \item En la vista Micro, examina los gráficos socioeconómicos y la distribución de suspensos para determinar si el problema requiere refuerzo curricular, becas, dotación tecnológica o programas de acompañamiento familiar.
    \item Con esta evidencia, fundamenta la propuesta de asignación de recursos ante los órganos competentes.
\end{enumerate}

Este flujo asegura que las decisiones se tomen sobre datos objetivos y no sobre impresiones subjetivas.